% Vietnamese translation for OI Public Library
% Source-PDF: 动态规划 DYNAMIC PROGRAMMING/1D1D动态规划优化初步.pdf
\documentclass[11pt]{article}
\usepackage{oipl-vn}

\title{Nhập môn tối ưu quy hoạch động 1D/1D}
\author{}
\date{}

\begin{document}
\maketitle

\origtitle{Introduction to 1D/1D Dynamic Programming Optimization}
\sourcepath{Dynamic Programming / 1D1D DP Optimization Primer.pdf}

\section*{Mở đầu}

Quy hoạch động 1D/1D là những phương trình có số trạng thái $O(n)$, và mỗi
trạng thái có $O(n)$ quyết định. Cách tính trực tiếp thường mất $O(n^2)$,
nhưng phần lớn các phương trình dạng này, nếu tổ chức hợp lý, có thể tối ưu
xuống $O(n\log n)$ hoặc thậm chí $O(n)$.

Bài viết này chỉ bàn về các mô hình và cách chuyển hóa kinh điển ở mức nhập
môn. Ta sẽ dùng hai ký hiệu cho hàm: $f(x)$ và $f[x]$. Giá trị viết bằng ngoặc
vuông được xem như có thể tính trước trước khi quy hoạch động bắt đầu, còn
giá trị viết bằng ngoặc tròn là biến phải tính trong quá trình quy hoạch động.
Một khi một giá trị đã được xác định, về sau nó không thay đổi nữa.

\section{Mô hình kinh điển thứ nhất}

Xét phương trình
\[
  f(x)=\min_{1\le i<x}\{f(i)+w[i,x]\}.
\]
Đây là mô hình rất quen thuộc. Một tính chất thường gặp là \term{đơn điệu
quyết định}: nếu $k(x)$ là quyết định tối ưu của trạng thái $x$, thì
\[
  i\le j \Longrightarrow k(i)\le k(j).
\]
Một điều kiện đủ kinh điển là bất đẳng thức tứ giác
\[
  w[i,j]+w[i+1,j+1]\le w[i+1,j]+w[i,j+1].
\]
Chứng minh tính chất này có thể tìm trong nhiều tài liệu về bất đẳng thức tứ
giác. Trong thực chiến, đôi khi cách đáng tin cậy hơn là viết thuật toán thô,
in bảng quyết định và quan sát; dù sao ta cũng cần đối chiếu với thuật toán
thô để kiểm thử.

Khi đã có đơn điệu quyết định, cách cài đặt thế nào cho hiệu quả? Ý tưởng
``từ trạng thái $x$ tìm quyết định tối ưu của nó'' không đủ mạnh: nếu chỉ
duyệt từ $k(x-1)$ trở đi thì chỉ giảm hằng số, còn nếu dừng ngay khi một quyết
định kém quyết định kế tiếp thì không đúng, vì đơn điệu quyết định không nói
rằng hàm $f(j)+w[j,x]$ đơn điệu theo $j$.

Ta đổi góc nhìn: với một trạng thái $f(j)$ đã tính, hãy xét nó có thể cập
nhật những trạng thái nào. Trong quá trình chạy, quyết định tạm thời của một
số trạng thái có thể chưa tối ưu, nhưng đến khi thuật toán kết thúc thì mọi
quyết định đều đúng.

Ban đầu chỉ có $f(1)$, nên mọi trạng thái chưa tính đều tạm lấy quyết định
$1$. Khi $f(2)$ được xác định, ta dùng quyết định $2$ cập nhật bảng quyết
định. Do tính đơn điệu, bảng chỉ có thể có dạng một đoạn quyết định $1$ rồi
đến một đoạn quyết định $2$. Vì vậy có thể dùng nhị phân để tìm điểm chuyển:
nếu tại một vị trí $x$, quyết định $2$ tốt hơn, thì mọi vị trí lớn hơn cũng
do quyết định $2$ thắng; nếu quyết định $1$ tốt hơn, thì mọi vị trí nhỏ hơn
cũng do quyết định $1$ thắng.

Tổng quát, dùng một ngăn xếp lưu các đoạn quyết định. Mỗi phần tử lưu một
quyết định và khoảng trạng thái $[b,e]$ mà quyết định này đang chiếm ưu thế.
Khi chèn một quyết định mới, quét từ cuối ngăn xếp:
\begin{enumerate}
  \item Nếu ở đầu đoạn $b$, quyết định mới đã tốt hơn quyết định cũ, thì bỏ
  toàn bộ đoạn cũ và gộp nó vào miền của quyết định mới.
  \item Nếu ở $b$ quyết định cũ vẫn tốt hơn, thì điểm chuyển nằm trong đoạn
  này; dùng nhị phân tìm điểm chuyển, đẩy quyết định mới vào ngăn xếp, rồi
  dừng.
\end{enumerate}
Mỗi quyết định bị loại tối đa một lần, nên phần quét ngăn xếp là khấu hao
$O(1)$; do còn nhị phân, độ phức tạp tổng là $O(n\log n)$.

\subsection{Ví dụ 1: Đóng gói đồ chơi}

Bài toán có $n$ đồ chơi, độ dài đồ chơi thứ $i$ là $c[i]$. Đồ chơi phải được
đóng theo thứ tự đã cho, và trong cùng một hộp, giữa hai đồ chơi kề nhau có
đúng một đơn vị khoảng cách. Nếu đóng các đồ chơi $i,\ldots,j$ vào cùng một
hộp, độ dài hộp là
\[
  \ell=j-i+\sum_{k=i}^{j}c[k].
\]
Chi phí của hộp độ dài $\ell$ là $(\ell-L)^2$, với $L$ cho trước. Cần đóng
tất cả đồ chơi với tổng chi phí nhỏ nhất.

Ta có phương trình 1D/1D tự nhiên
\[
  f(x)=\min_{1\le i<x}\{f(i)+w[i+1,x]\},
  \qquad
  w[i,j]=\left(j-i+\sum_{k=i}^{j}c[k]-L\right)^2.
\]
Hàm $w$ thỏa tính đơn điệu quyết định, nên có thể dùng phương pháp trên để
giải trong $O(n\log n)$.

\subsection{Ví dụ 2: Mua đất}

Bài USACO Gold tháng 3 năm 2008 cho $N$ mảnh đất hình chữ nhật. Một lần mua
một nhóm đất có giá bằng chiều dài lớn nhất nhân với chiều rộng lớn nhất
trong nhóm. Cần chia nhóm để tổng giá nhỏ nhất.

Sắp các mảnh theo chiều dài giảm dần. Khi quét theo thứ tự này, chiều dài
của mảnh hiện tại luôn không lớn hơn các mảnh trước đó, nên chỉ cần quan tâm
đến chiều rộng. Nếu một điểm $(x,y)$ bị một điểm khác $(x',y')$ chi phối, tức
$x'\ge x$ và $y'\ge y$, thì điểm đó có thể bỏ qua vì nó luôn được một nhóm
khác bao phủ. Sau một lần quét tuyến tính, các điểm còn lại có chiều dài và
chiều rộng biến thiên ngược chiều nhau.

Với dãy đã rút gọn, phương trình trở thành
\[
  f(n)=\min_{0\le k<n}\{f(k)+x[n]\cdot y[k+1]\}.
\]
Đây là dạng có đơn điệu quyết định, nên cũng giải được bằng phương pháp
ngăn xếp ở trên trong $O(n\log n)$.

Hai ví dụ này đều là ứng dụng trực tiếp của đơn điệu quyết định. Ví dụ mua
đất nhắc ta rằng trước khi cố tối ưu phương trình, đôi khi cần loại những
quyết định chắc chắn vô dụng và sắp xếp lại dữ liệu; nếu giữ nguyên dữ liệu
ban đầu thì phương trình có thể không còn đơn điệu.

\section{Mô hình kinh điển thứ hai}

Xét một lớp đặc biệt của $w$:
\[
  \forall i\le j<k,\qquad w[i,j]+w[j,k]=w[i,k].
\]
Khi đó tồn tại một hàm một biến $w'$ sao cho
\[
  w[i,j]=w'[j]-w'[i].
\]
Nếu quyết định của trạng thái $x$ là $k$, thì
\[
\begin{aligned}
  f(x)
    &=f(k)+w[k,x] \\
    &=f(k)+w'[x]-w'[k] \\
    &=g(k)+w'[x]-w'[1],
\end{aligned}
\]
trong đó $g(k)=f(k)-w[1,k]$. Phần phụ thuộc vào $x$ là hằng số đối với mọi
quyết định $k$, nên nếu không có ràng buộc gì thêm, ta chỉ cần duy trì
$\min g(k)$ bằng một biến.

Khi có cận dưới đơn điệu cho quyết định, ta nhận được mô hình thứ hai:
\[
  f(x)=\min_{b[x]\le k<x}\{g(k)\}+w[x],
  \qquad b[x]\ \text{không giảm theo }x.
\]
Nếu có hai quyết định $j\le k$ và $g(k)\le g(j)$, thì $j$ vô dụng: mỗi khi
$j$ hợp lệ, $k$ cũng hợp lệ, và $k$ tốt không kém $j$. Vì vậy nếu xếp quyết
định theo chỉ số, các giá trị $g$ trong hàng đợi phải tăng. Ta dùng hàng đợi
đơn điệu:
\begin{enumerate}
  \item Loại các phần tử đầu hàng đợi cho đến khi phần tử đầu nằm trong phạm
  vi quyết định hợp lệ.
  \item Phần tử đầu là quyết định tối ưu của trạng thái hiện tại.
  \item Tính $g(x)$ và chèn vào cuối hàng đợi; trong lúc chèn, loại các phần
  tử cuối có giá trị $g$ không tốt hơn phần tử mới.
\end{enumerate}
Mỗi phần tử vào và ra hàng đợi tối đa một lần, nên thời gian khấu hao là
$O(1)$ cho mỗi trạng thái.

\subsection{Ví dụ 3: The Sound of Silence}

Bài Baltic OI 2007 cho một dãy $N$ số. Một đoạn liên tiếp độ dài $M$ là hợp
lệ nếu hiệu giữa giá trị lớn nhất và nhỏ nhất trong đoạn không vượt quá hằng
số $C$. Cần in mọi vị trí bắt đầu của đoạn hợp lệ.

Bài toán tách thành hai phần giống nhau: tìm cực đại trượt và cực tiểu trượt
trên mọi cửa sổ độ dài $M$. Chẳng hạn, nếu $g[x]$ là giá trị phần tử thứ
$x$, còn $f(x)$ là giá trị nhỏ nhất trong cửa sổ kết thúc tại $x$, thì
\[
  f(x)=\min_{x-M+1\le i\le x} g[i].
\]
Đây chính là mô hình thứ hai, nên hàng đợi đơn điệu cho lời giải $O(n)$.

\subsection{Ví dụ 4: Islands}

Bài IOI 2008 yêu cầu, với một đồ thị vô hướng có trọng số gồm $N$ đỉnh và
đúng $N$ cạnh, tìm đường đi đơn dài nhất trong mỗi thành phần liên thông. Mỗi
thành phần có cấu trúc một chu trình cơ sở kèm các cây mọc ra từ các đỉnh
trên chu trình.

Đường đi dài nhất hoặc nằm hoàn toàn trong một cây mọc ra, hoặc đi qua một
đoạn trên chu trình và nối hai cây mọc ra. Trường hợp thứ nhất xử lý bằng DP
trên cây. Với trường hợp thứ hai, phá chu trình thành dây rồi nhân đôi dây
để mọi đoạn trên chu trình đều tương ứng với một đoạn liên tiếp. Gọi $g[x]$
là đường đi đi xuống dài nhất trong cây mọc ra tại đỉnh $x$. Khi lấy $x$ làm
đầu phải của đoạn, ta có
\[
  f(x)=\max_{x-n+1\le i<x}\{g[i]+g[x]+\operatorname{dis}(i,x)\}.
\]
Khoảng cách trên dây thỏa tính cộng được như $w[i,j]+w[j,k]=w[i,k]$, nên sau
biến đổi, bài toán rơi vào mô hình hàng đợi đơn điệu và xử lý tổng cộng trong
$O(n)$.

\subsection{Ví dụ 5: Balo hữu hạn}

Trong balo hữu hạn, vật thứ $k$ có giá trị $p[k]$, trọng lượng $w[k]$, và
được chọn tối đa $m[k]$ lần. Nếu dùng $f(x)$ là trọng lượng nhỏ nhất để đạt
được tổng giá trị $x$, khi xét vật $k$ ta có
\[
  f(x)=\min_{1\le i\le m[k]}
       \{f(x-i\,p[k])+i\,w[k]\}.
\]
Quyết định trong phương trình này không liên tiếp theo $x$, vì các trạng thái
hợp lệ cách nhau $p[k]$. Chia các trạng thái theo lớp đồng dư modulo $p[k]$,
trong mỗi lớp quyết định trở thành một đoạn liên tiếp. Nếu gọi hàm mới là
$h$, phương trình có dạng
\[
  h(x)=\min_{x-m[k]\le i<x}
       \{h(i)+w[k](x-i)\},
\]
trong đó phần $w[k](x-i)$ là một hàm cộng được. Vì vậy nó chuyển được về mô
hình thứ hai. Đây chính là cách hàng đợi đơn điệu trong bài balo hữu hạn.

Nếu $m[k]=1$, mỗi trạng thái chỉ có $O(1)$ quyết định. Nếu $m[k]=+\infty$,
cận dưới biến mất và chỉ cần duy trì một giá trị tốt nhất. Nhìn từ mô hình
thứ hai, balo 0--1, balo hữu hạn và balo vô hạn được đặt trong cùng một khung
nhìn.

\section{Mô hình kinh điển thứ ba}

Quay lại phương trình mua đất đã rút gọn:
\[
  f(n)=\min_{0\le k<n}\{f(k)+x[n]\,y[k+1]\}.
\]
Ở đây $f(k)$ là biến, $y[k+1]$ là giá trị đã biết sau khi trạng thái $k$ được
xác định, còn $x[n]$ là hệ số chỉ phụ thuộc vào trạng thái đang tính. Nếu ta
xem mỗi quyết định $k$ là một điểm $(X,Y)=(y[k+1],f(k))$ trên mặt phẳng, thì
việc tính $f(n)$ là tối thiểu hóa một biểu thức tuyến tính
\[
  P=aX+bY.
\]
Đường thẳng mức có dạng
\[
  Y=-\frac{a}{b}X+\frac{P}{b}.
\]
Vì vậy bài toán trở thành chọn điểm trên mặt phẳng sao cho một đường thẳng có
hệ số góc xác định đạt tung độ cắt nhỏ nhất. Các điểm tối ưu luôn nằm trên
bao lồi của tập điểm quyết định.

Mô hình tổng quát là
\[
  f(n)=\min_{1\le i<n}\{a[n]\,x(i)+b[n]\,y(i)\},
\]
trong đó cặp $(x(i),y(i))$ được xác định duy nhất trong $O(1)$ từ trạng thái
$f(i)$. Ta có hai tình huống chính.

\paragraph{Trường hợp 1: hệ số góc và hoành độ điểm quyết định có tính đơn
điệu.}
Khi hệ số góc thay đổi đơn điệu, điểm quyết định tối ưu cũng di chuyển đơn
điệu trên bao lồi. Ta duy trì bao lồi bằng hàng đợi đơn điệu:
\begin{enumerate}
  \item Dời con trỏ quyết định ở đầu hàng đợi cho đến điểm tốt nhất hiện tại.
  \item Dùng điểm đó để tính trạng thái.
  \item Chèn điểm mới vào cuối hàng đợi, dùng thao tác kiểu Graham scan để
  loại các điểm làm mất tính lồi.
\end{enumerate}
Mỗi điểm bị loại tối đa một lần, nên độ phức tạp là $O(n)$.

\paragraph{Trường hợp 2: không có ràng buộc đơn điệu.}
Điểm tối ưu vẫn nằm trên bao lồi, nhưng cần tìm bằng nhị phân theo hệ số góc.
Việc chèn điểm mới cũng khó hơn: cần tìm vị trí theo hoành độ, rồi duy trì
tính lồi ở hai phía. Một cách hoàn chỉnh là dùng cây cân bằng theo hoành độ;
sau khi chèn, đưa điểm mới lên gốc và cắt bỏ những đoạn không còn thuộc bao
lồi. Độ phức tạp lý thuyết là $O(n\log n)$, nhưng cài đặt phức tạp.

\subsection{Ví dụ 6: Thuật toán tuyến tính cho đóng gói đồ chơi}

Với bài đóng gói đồ chơi, đặt
\[
  \operatorname{sum}[x]=\sum_{i=1}^{x}c[i].
\]
Phương trình có thể viết lại:
\[
  f(x)=
  \min_{1\le i<x}
  \left\{
    f(i)+
    \bigl(x-i-1-L+\operatorname{sum}[x]-\operatorname{sum}[i]\bigr)^2
  \right\}.
\]
Đặt
\[
  a[x]=\operatorname{sum}[x]+x-L-1,
  \qquad
  b[i]=\operatorname{sum}[i]+i.
\]
Khi đó
\[
\begin{aligned}
  f(x)
  &=\min_i \{f(i)+(a[x]-b[i])^2\}\\
  &=a[x]^2+\min_i\{f(i)+b[i]^2-2a[x]b[i]\}.
\end{aligned}
\]
Xem mỗi quyết định $i$ là điểm
\[
  X(i)=b[i],\qquad Y(i)=f(i)+b[i]^2.
\]
Ta cần tối thiểu hóa $Y-2a[x]X$. Vì $a[x]$ tăng theo $x$ và $X(i)$ tăng theo
$i$, hệ số góc và hoành độ đều đơn điệu; dùng hàng đợi bao lồi sẽ cho thuật
toán $O(n)$.

\subsection{Ví dụ 7: Đổi tiền}

Bài NOI 2007 có ba loại tiền: nhân dân tệ, chứng khoán A và chứng khoán B.
Mỗi ngày giá của A và B thay đổi. Trong ngày $D$, ta có thể bán toàn bộ A và
B lấy nhân dân tệ, mua A và B theo một tỉ lệ cố định bằng toàn bộ số nhân
dân tệ đang có, hoặc không làm gì. Ban đầu có một lượng nhân dân tệ, cần tối
đa hóa số nhân dân tệ ở cuối ngày cuối cùng.

Gọi $f(n)$ là số nhân dân tệ nhiều nhất có thể có đến ngày $n$, và $(x(i),y(i))$
là lượng A và B nhận được nếu dùng số tiền tối ưu của ngày $i$ để mua theo tỉ
lệ đã cho. Nếu $a[n]$ và $b[n]$ là giá A và B ở ngày $n$, ta có
\[
  f(n)=\max_{1\le i<n}\{a[n]\,x(i)+b[n]\,y(i)\}.
\]
Đây là mô hình bao lồi trong trường hợp không có ràng buộc đơn điệu rõ ràng,
nên cách chuẩn là dùng cây cân bằng để duy trì bao lồi động, đạt
$O(n\log n)$. Với dữ liệu yếu, một số cài đặt dùng con trỏ di chuyển cũng có
thể qua, nhưng đó không phải phân tích tổng quát.

\section{Tổng kết}

Bài viết đã tập trung vào ba mô hình 1D/1D:
\[
  f(x)=\min_{1\le i<x}\{f(i)+w[i,x]\},
\]
trong đó $w$ thỏa đơn điệu quyết định;
\[
  f(x)=\min_{b[x]\le i<x}\{g(i)\}+w[x],
\]
trong đó $b[x]$ đơn điệu tăng; và
\[
  f(n)=\min_{1\le i<n}\{a[n]\,x(i)+b[n]\,y(i)\},
\]
trong đó $(x(i),y(i))$ được xác định từ $f(i)$ trong thời gian hằng số.
Ba mô hình này đều có thể giải ít nhất trong $O(n\log n)$, và trong nhiều
trường hợp đạt $O(n)$. Chúng cũng không tách rời nhau: nhiều ví dụ trong bài
có thể được nhìn dưới nhiều mô hình khác nhau sau khi biến đổi đại số.

Đây mới chỉ là thảo luận nhập môn. Một số hướng còn đáng nghiên cứu:
\begin{enumerate}
  \item Trong mô hình thứ nhất, liệu có thuật toán $O(n)$ tổng quát cho mọi
  hàm $w$ thỏa đơn điệu quyết định hay không? Bài viết chỉ nêu thuật toán
  $O(n\log n)$ tổng quát và một số trường hợp đặc biệt $O(n)$.
  \item Trong hai mô hình đầu, ta có thể thêm cận dưới đơn điệu $b[x]$ mà
  không ảnh hưởng đến độ phức tạp. Với mô hình bao lồi, do thứ tự chọn quyết
  định phụ thuộc vào mặt phẳng các điểm $(x(i),y(i))$ chứ không phụ thuộc đơn
  thuần vào thứ tự trạng thái, việc thêm $b[x]$ có còn tối ưu được hay không
  là một câu hỏi.
  \item Các phương trình đã xét đều có tập quyết định liên tiếp. Nhiều bài
  1D/1D kinh điển, chẳng hạn LIS, chỉ cho phép những quyết định có khóa nhỏ
  hơn khóa hiện tại, nên không trực tiếp dùng được các kỹ thuật trên. Tối ưu
  các phương trình 1D/1D tổng quát hơn vẫn là vấn đề đáng khảo sát.
  \item 1D/1D cũng có giới hạn của nó: số trạng thái vẫn là một chiều. Trong
  thực tế còn nhiều phương trình 2D/0D như LCS, đua ngựa Tian Ji, sửa chữa
  công trình hay các bài tương tự. Hệ phương pháp tối ưu cho những dạng này
  chưa phong phú bằng 1D/1D.
\end{enumerate}

\section*{Bài tập}

\begin{enumerate}
  \item Cài đặt hai cách giải cho bài đóng gói đồ chơi và bài mua đất.
  \item Bài \emph{Vũ điệu lộng lẫy} của NOI 2005: cho một bàn cờ $N\times M$
  có ô trống và vật cản. Từ một điểm đầu, trong từng khoảng thời gian liên
  tiếp, hướng trượt được cố định; mỗi thời điểm có thể đứng yên hoặc trượt
  một bước theo hướng đó, miễn là không đâm vào vật cản. Hãy cho thuật toán
  $O(NMK)$. Gợi ý: bài này có thể đối ứng trực tiếp với balo hữu hạn.
  \item \emph{Cactus Reloaded}: trong đồ thị xương rồng liên thông, mỗi cạnh
  thuộc nhiều nhất một chu trình. Cần tìm khoảng cách ngắn nhất lớn nhất giữa
  hai đỉnh bất kỳ. Hãy cho thuật toán $O(N+M)$.
  \item Cài đặt bài đổi tiền bằng cây cân bằng để duy trì bao lồi.
\end{enumerate}

\end{document}
